\documentclass[openany]{book}

\usepackage[margin=1in]{geometry}
\usepackage{amsmath,amsfonts,amsthm, amssymb}
\usepackage{yhmath}
\usepackage{mathrsfs}
\usepackage{mathtools}
\usepackage{xcolor}
\usepackage{graphicx}
\usepackage{comment}
\usepackage{tikz-cd}
\usepackage{quiver}
\renewcommand{\familydefault}{ppl}
\newcommand{\R}{\mathbb{R}}
\newcommand{\E}{\mathbb{E}}
\newcommand{\Z}{\mathbb{Z}}
\newcommand{\CC}{\mathbb{C}}
\newcommand{\F}{\mathbb{F}}
\newcommand{\la}{\langle}
\newcommand{\ra}{\rangle}
\newcommand{\colim}{\text{colim}}
\DeclareMathOperator{\im}{im}
\let\oldemptyset\emptyset
\let\emptyset\varnothing
\newcommand{\tor}{\text{Tor}}
\newcommand{\id}{\text{id}}
\newcommand{\ext}{\text{Ext}}
\newcommand{\ptop}{\text{PTop}}
\newcommand{\pt}{\text{pt}}
\newcommand{\ach}{\text{Ach}}
\newcommand{\Q}{\mathbb{Q}}


\usepackage{thmtools,thm-restate}

% Fixing mdframed skip below
% See https://tex.stackexchange.com/a/292090/143086
\usepackage[framemethod=TikZ]{mdframed}
\usepackage{xpatch}
\makeatletter
\xpatchcmd{\endmdframed}
	{\aftergroup\endmdf@trivlist\color@endgroup}
	{\endmdf@trivlist\color@endgroup\@doendpe}
	{}{}
\makeatother

\definecolor{huilightpink}{HTML}{fff2fe}
\definecolor{huidarkpink}{HTML}{d955b7}
\declaretheoremstyle[
	mdframed={
		backgroundcolor=huilightpink,
		linecolor=huidarkpink,
		rightline=false,
		topline=false,
		bottomline=false,
		linewidth=2pt,
		innertopmargin=5pt,
		innerbottommargin=8pt,
		innerleftmargin=8pt,
		leftmargin=-2pt,
		skipbelow=2pt,
		nobreak
	},
	headfont=\normalfont\bfseries\color{huidarkpink}
]{huipinkbox}
\declaretheorem[style=huipinkbox,name=Theorem,within=chapter]{thm}
\declaretheorem[style=huipinkbox,name=Theorem,sibling=thm]{theorem}




\begin{comment}
\definecolor{huilightyellow}{HTML}{fff5d6}
\definecolor{huidarkyellow}{HTML}{fcad03}
\declaretheoremstyle[
	mdframed={
		backgroundcolor=huilightyellow,
		linecolor=huidarkyellow,
		rightline=false,
		topline=false,
		bottomline=false,
		linewidth=2pt,
		innertopmargin=5pt,
		innerbottommargin=8pt,
		innerleftmargin=8pt,
		leftmargin=-2pt,
		skipbelow=2pt,
		nobreak
	},
	headfont=\normalfont\bfseries\color{huidarkyellow}
]{huiyellowbox}
\declaretheorem[style=huiyellowbox,name=Proposition,within=chapter]{prop}
\end{comment}



\definecolor{huilightpurple}{HTML}{faf2ff}
\definecolor{huidarkpurple}{HTML}{912ed9}
\declaretheoremstyle[
	mdframed={
		backgroundcolor=huilightpurple,
		linecolor=huidarkpurple,
		rightline=false,
		topline=false,
		bottomline=false,
		linewidth=2pt,
		innertopmargin=5pt,
		innerbottommargin=8pt,
		innerleftmargin=8pt,
		leftmargin=-2pt,
		skipbelow=2pt,
		nobreak
	},
	headfont=\normalfont\bfseries\color{huidarkpurple}
]{huipurplebox}
\declaretheorem[style=huipurplebox,name=Proposition,within=chapter]{prop}



% \definecolor{huilightpurple}{HTML}{faf2ff}
% \definecolor{huidarkpurple}{HTML}{912ed9}
% \declaretheoremstyle[
% 	mdframed={
% 		backgroundcolor=huilightpurple,
% 		linecolor=huidarkpurple,
% 		rightline=false,
% 		topline=false,
% 		bottomline=false,
% 		linewidth=2pt,
% 		innertopmargin=5pt,
% 		innerbottommargin=8pt,
% 		innerleftmargin=8pt,
% 		leftmargin=-2pt,
% 		skipbelow=2pt,
% 		nobreak
% 	},
% 	headfont=\normalfont\bfseries\color{huidarkpurple}
% ]{huipurplebox}
\declaretheorem[style=huipurplebox,name=Lemma,within=chapter]{lem}


\definecolor{lightpink}{HTML}{f0f6fc}
\definecolor{darkpink}{HTML}{2c72b8}
\declaretheoremstyle[
	mdframed={
		backgroundcolor=lightpink,
		linecolor=darkpink,
		rightline=false,
		topline=false,
		bottomline=false,
		linewidth=2pt,
		innertopmargin=5pt,
		innerbottommargin=8pt,
		innerleftmargin=8pt,
		leftmargin=-2pt,
		skipbelow=2pt,
		nobreak
	},
	headfont=\normalfont\bfseries\color{darkpink}
]{pinkbox}
\declaretheorem[style=pinkbox,name=Definition,within=chapter]{defn}


\definecolor{huilightblue}{HTML}{edf9ff}
\definecolor{huidarkblue}{HTML}{4b79db}
\declaretheoremstyle[
	mdframed={
		backgroundcolor=huilightblue,
		linecolor=huidarkblue,
		rightline=false,
		topline=false,
		bottomline=false,
		linewidth=2pt,
		innertopmargin=5pt,
		innerbottommargin=8pt,
		innerleftmargin=8pt,
		leftmargin=-2pt,
		skipbelow=2pt,
		nobreak
	},
	headfont=\normalfont\bfseries\color{huidarkblue}
]{huiblueblox}
\declaretheorem[style=huiblueblox,name=Example,within=chapter]{example}



% \definecolor{huilightblue}{HTML}{edf9ff}
% \definecolor{huidarkblue}{HTML}{4b79db}
% \declaretheoremstyle[
% 	mdframed={
% 		backgroundcolor=huilightblue,
% 		linecolor=huidarkblue,
% 		rightline=false,
% 		topline=false,
% 		bottomline=false,
% 		linewidth=2pt,
% 		innertopmargin=5pt,
% 		innerbottommargin=8pt,
% 		innerleftmargin=8pt,
% 		leftmargin=-2pt,
% 		skipbelow=2pt,
% 		nobreak
% 	},
% 	headfont=\normalfont\bfseries\color{huidarkblue}
% ]{huiblueblox}
% \declaretheorem[style=huiblueblox,name=Example,within=chapter]{example}

% \declaretheoremstyle[
% 	mdframed={
% 		backgroundcolor=huilightblue,
% 		linecolor=huidarkblue,
% 		rightline=false,
% 		topline=false,
% 		bottomline=false,
% 		linewidth=2pt,
% 		innertopmargin=5pt,
% 		innerbottommargin=8pt,
% 		innerleftmargin=8pt,
% 		leftmargin=-2pt,
% 		skipbelow=2pt,
% 		nobreak
% 	},
% 	headfont=\normalfont\bfseries\color{huidarkblue}
% ]{huiblueblox}
\declaretheorem[style=huiblueblox,name=Problem,within=chapter]{prob}



% \declaretheoremstyle[
% 	mdframed={
% 		backgroundcolor=huilightblue,
% 		linecolor=huidarkblue,
% 		rightline=false,
% 		topline=false,
% 		bottomline=false,
% 		linewidth=2pt,
% 		innertopmargin=5pt,
% 		innerbottommargin=8pt,
% 		innerleftmargin=8pt,
% 		leftmargin=-2pt,
% 		skipbelow=2pt,
% 		nobreak
% 	},
% 	headfont=\normalfont\bfseries\color{huidarkblue}
% ]{huiblueblox}
\declaretheorem[style=huiblueblox,name=Exercise,within=chapter]{exer}
\declaretheorem[style=huipinkbox, name=Corollary, within=chapter]{cor}





















\newcommand{\nirwarnsymbol}{%
	\begin{tikzpicture}[baseline=(x.base)]
		\draw[rounded corners=.01em] (-.05em,-1.07em)rectangle(.05em,.78em);
		\draw[fill=white,rounded corners=1.3] (0,.75em)--(.75em,0)--(0,-.75em)--(-.75em,0)--cycle;
		\draw[line width=0.2mm, line cap=round](-.4em,-1.07em)--(.4em,-1.07em);
		\node(x) at (0,0em) {};
		% Thank you https://tex.stackexchange.com/a/262510
		\draw[
			line cap=but,
			line join=round,
			x=.5em,
			line width=0.5mm,
			y=1*(height("Z")-\pgflinewidth)*(1-sin(10)),
			rotate=-10,
			rounded corners=1.5pt,
		](-0.57, 0.57) -- (0.57, 0.57) -- (-0.57, -0.57) -- (0.57, -0.57);
	\end{tikzpicture}%
}

%%%%%%%%%%%%%%%%%%%%%%%%%%%%%%%%%%%%%%%%%%%% MARGINS
\usepackage{marginnote}
% Thank you https://tex.stackexchange.com/a/472882
% Makes marginnotes always appear on the left, apparently
%
\makeatletter
\long\def\@mn@@@marginnote[#1]#2[#3]{%
	\begingroup
		\ifmmode\mn@strut\let\@tempa\mn@vadjust\else
			\if@inlabel\leavevmode\fi
			\ifhmode\mn@strut\let\@tempa\mn@vadjust\else\let\@tempa\mn@vlap\fi
		\fi
		\@tempa{%
			\vbox to\z@{%
				\vss
				\@mn@margintest
				\if@reversemargin\if@tempswa
						\@tempswafalse
					\else
						\@tempswatrue
				\fi\fi

					\llap{%
						\vbox to\z@{\kern\marginnotevadjust\kern #3
							\vbox to\z@{%
								\hsize\marginparwidth
								\linewidth\hsize
								\kern-\parskip
								%\mn@parboxrestore
								\marginfont\raggedleftmarginnote\strut\hspace{\z@}%
								\ignorespaces#1\endgraf
								\vss
							}%
							\vss
						}%
						\if@mn@verbose
							\PackageInfo{marginnote}{xpos seems to be \@mn@currxpos}%
						\fi
						\begingroup
							\ifx\@mn@currxpos\relax\else\ifx\@mn@currpos\@empty\else
									\kern\@mn@currxpos
							\fi\fi
							\ifx\@mn@currpage\relax
								\let\@mn@currpage\@ne
							\fi
							\if@twoside\ifodd\@mn@currpage\relax
									\kern-\oddsidemargin
								\else
									\kern-\evensidemargin
								\fi
							\else
								\kern-\oddsidemargin
							\fi
							\kern-1in
						\endgroup
						\kern\marginparsep
					}%
			}%
		}%
	\endgroup
}
\makeatother
%
% Mostly for todonotes
\renewcommand{\marginpar}{\marginnote}
%%%%%%%%%%%%%%%%%%%%%%%%%%%%%%%%%%%%%%%%%%%% /MARGINS

\definecolor{nirlightred}{RGB}{250, 220, 220}
\definecolor{nirdarkred}{HTML}{f40000}
\declaretheoremstyle[
	mdframed={
		backgroundcolor=nirlightred,
		linecolor=nirdarkred,
		rightline=false,
		topline=false,
		bottomline=false,
		linewidth=2pt,
		innertopmargin=5pt,
		innerbottommargin=8pt,
		innerleftmargin=8pt,
		leftmargin=-2pt,
		skipbelow=2pt,
		nobreak
	},
	headfont=\normalfont\bfseries\color{nirdarkred}
]{nirredbox}

% \makeatletter
% \declaretheorem[
% 	style=nirredbox,
% 	name=Warning,
% 	sibling=thm,
% 	% without \leavevmode, the first item in a list gets misformatted
% 	postheadhook={\leavevmode\marginnote{\nirwarnsymbol}[-3pt]%
% 	\ifthmt@thisistheone% restatable makes alignment weird
% 		\hspace{-2.2pt}%
% 	\fi}
% ]{warn}
% \makeatother

\newcommand{\nirideasymbol}{%
	\begin{tikzpicture}[baseline=(x.base)]
		\draw[rounded corners=.01em] (-.05em,-1.07em)rectangle(.05em,.78em);
		\draw[fill=white,rounded corners=1.3] (0,.75em)--(.75em,0)--(0,-.75em)--(-.75em,0)--cycle;
		\draw[line width=0.2mm, line cap=round](-.4em,-1.07em)--(.4em,-1.07em);
		\node(x) at (0,0em) {};
		\node at (0,0em) {{\textbf{!}}};
	\end{tikzpicture}%
}
\renewcommand{\nirwarnsymbol}{%
	\begin{tikzpicture}[baseline=(x.base)]
		\draw[rounded corners=.01em] (-.05em,-1.07em)rectangle(.05em,.78em);
		\draw[fill=white,rounded corners=1.3] (0,.75em)--(.75em,0)--(0,-.75em)--(-.75em,0)--cycle;
		\draw[line width=0.2mm, line cap=round](-.4em,-1.07em)--(.4em,-1.07em);
		\node(x) at (0,0em) {};
		% Thank you https://tex.stackexchange.com/a/262510
		\draw[
			line cap=but,
			line join=round,
			x=.5em,
			line width=0.5mm,
			y=1*(height("Z")-\pgflinewidth)*(1-sin(10)),
			rotate=-10,
			rounded corners=1.5pt,
		](-0.57, 0.57) -- (0.57, 0.57) -- (-0.57, -0.57) -- (0.57, -0.57);
	\end{tikzpicture}%
}
\makeatletter
\declaretheorem[
	style=nirredbox,
	name=Idea,
	sibling=thm,
	% without \leavevmode, the first item in a list gets misformatted
	postheadhook={\leavevmode\marginnote{\nirideasymbol}[-3pt]%
	\ifthmt@thisistheone% restatable makes alignment weird
		\hspace{-2.2pt}%
	\fi}
]{idea}

\declaretheorem[
	style=nirredbox,
	name=Warning,
	sibling=thm,
	% without \leavevmode, the first item in a list gets misformatted
	postheadhook={\leavevmode\marginnote{\nirwarnsymbol}[-3pt]%
	\ifthmt@thisistheone% restatable makes alignment weird
		\hspace{-2.2pt}%
	\fi}
]{warn}
\makeatother

\title{Algebra Qualifying Exam Solutions}

\date{\today}
\author{Hui Sun}


\begin{document}

\maketitle

\tableofcontents
\newpage



\chapter{Fall 2016}



\begin{prob}
    Determine \(\text{Aut}(S_{3})\).
\end{prob}
\begin{proof}
    $\sigma\in\text{Aut}(S_3)$ is determined by where $(12)$ and $(123)$ are sent to. There are $6$ options in total and all of them are homomorphisms (conjugation). It is easy to check that this group is not commutative, i.e., 
    \begin{equation*}
        \text{Aut}(S_3)\cong S_3
    \end{equation*}
\end{proof}

\begin{prob}
    A group \(G\) is a semidirect product of subgroups \(N,H\subset G\) if \(N\) is normal and every element of \(G\) has a unique presentation \(nh\), \(n\in N\), \(h\in H\). Find all semidirect products (up to isomorphism) of \(N=\mathbb{Z}/11\mathbb{Z}\), \(H=\mathbb{Z}/5\mathbb{Z}\).
\end{prob}
\begin{proof}
    Let $G=N\rtimes_\theta H$, where 
    \begin{equation*}
        \theta: \Z/5\Z\to \text{Aut}(\Z/11\Z)\cong\Z/10\Z
    \end{equation*}
    such that
    \begin{equation*}
        5\theta(1)\equiv 0\mod 10
    \end{equation*}
    Thus $\theta(1)$ could be $0, 2,4,6,8$. When $\theta(1)=0$, this gives the abelian group 
    \begin{equation*}
        G\cong\frac{\Z}{5\Z}\times\frac{\Z}{11\Z}
    \end{equation*}
    We claim that all nontrivial $\theta$ give rise to the same semidirect product, namely, the following diagram commutes
    \[\begin{tikzcd}
        {\Z/5\Z} & {\Z/10\Z} \\
        {\Z/5\Z} & {\Z/10\Z}
        \arrow["{\theta'}", from=1-1, to=1-2]
        \arrow["m"', from=1-1, to=2-1]
        \arrow["{\text{id}}", from=1-2, to=2-2]
        \arrow["\theta", from=2-1, to=2-2]
    \end{tikzcd}\]
    for $\theta: 1\mapsto 2$ and any $\theta': 1\mapsto 4, 6, 8$, by taking $m$ to be the multiplication map by $2,3,4$ respectively. Hence we see 
    \begin{equation*}
        \theta(h)(g)=g^{2^{2h}}
    \end{equation*}
    by observing 
    \[\begin{tikzcd}
        {\Z/5\Z} & {\Z/10\Z} & {\left(\Z/11\Z\right)^\times} & {\text{Aut}(\Z/11\Z)}
        \arrow["2", from=1-1, to=1-2]
        \arrow["{2^2}", from=1-2, to=1-3]
        \arrow["{2^2\cdot(-)}"', from=1-3, to=1-4]
    \end{tikzcd}\]
    In other words,
    \begin{equation*}
        G=\left\la g,h: g^11=h^5=1, hgh^{-1}=g^{2^{2h}}\right\ra
    \end{equation*}
\end{proof}

\begin{prob}
    Let \(F\) be a finite field of order \(2^{n}\). Here \(n>0\). Determine all values of \(n\) such that the polynomial \(x^{2}-x+1\) is irreducible in \(F[x]\).
\end{prob}
\begin{proof}
    We know that $x^2-x+1$ is irreducible over $\F_2$, namely, it has no roots in $\F_2$. Since there is only one field of order $4$, we must have 
    \begin{equation*}
        \F_4\cong\frac{\F_2}{(x^2-x+1)}
    \end{equation*}
    Clearly $x^2-x+1$ is not irreducible over $\F_4$. For any $\F_{2^n}$, we know $(x^2-x+1)$ is irreducible if and only if $\F_4$ does not embed into $\F{2^n}$, i.e., $2\nmid n$. This shows that when $n$ is odd, the polynomial $x^2-x+1$ is irreducible over $\F_{2^n}$.
\end{proof}

\begin{prob}
    \begin{itemize}
        \item[(1)] Determine the Galois group of \(x^{4}-4x^{2}-2\) over \(\mathbb{Q}\).
        \item[(2)] Let \(G\) be a group of order 8 such that \(G\) is the Galois group of a polynomial of degree 4 over \(\mathbb{Q}\). Show that \(G\) is isomorphic to the Galois group in part (1).
    \end{itemize}
\end{prob}
\begin{proof}
    \begin{itemize}
        \item[(1)] The roots of this polynomial is $\pm\sqrt{2\pm\sqrt{6}}$, and notice that 
        \begin{equation*}
            \sqrt{2}i=\sqrt{2+\sqrt{6}}\sqrt{2-\sqrt{6}}
        \end{equation*}
        This gives the splitting field (Galois extension) of this polynomila as 
        \begin{equation*}
            \Q\left(\sqrt{2+\sqrt{6}}, \sqrt{2}i\right)
        \end{equation*}
        We see that 
        \begin{equation*}
            \Q\left(\sqrt{2+\sqrt{6}}\right)\cap\Q(\sqrt{2}i)=\emptyset
        \end{equation*}
        because the first is contained in $\R$ and the second is not. We must have 
        \begin{equation*}
            \left[ \Q\left(\sqrt{2+\sqrt{6}}, \sqrt{2}i\right)/\Q\right]=8
        \end{equation*}
        By part b, we see $\text{Gal}\cong D_8$.
        \item[(2)] Any Galois group of a polynomial with $4$ roots in the splitting field embeds into $S_4$, and we notice that $|G|=2^3, |S_4|=2^3\cdot 3$, i.e., $G$ is a Sylow $2$-subgroup of $S_4$, and all Sylow $2$-subgroups are conjugate/isomorphic of one another, hence 
        \begin{equation*}
            \text{Gal}\cong D_8
        \end{equation*}
    \end{itemize}
\end{proof}

\begin{prob}
    Let A be a linear transformation of a finite dimensional vector space over a field of characteristic \(\neq 2\).
    \begin{itemize}
        \item[(1)] Define the wedge product linear transformation \(\wedge^{2}A=A\wedge A\).
        \item[(2)] Prove that
        \[tr(\wedge^{2}A)=\frac{1}{2}(tr(A)^{2}-tr(A^{2})).\]
    \end{itemize}
\end{prob}
\begin{proof}
    (Recall we have analogous results for $A\otimes A$).
    \begin{itemize}
        \item[(1)] The wedge product $A\wedge A$ is defined on the wedge product of vector spaces $V\wedge V$, so we first define the vector space: let $\{v_1,\dots, v_n\}$ be the basis of $V$, then $\{v_i\wedge v_j\}$ where $i<j$ forms a basis of $V\wedge V$, satisfying:
        \begin{enumerate}
            \item $v_i\wedge v_j=-v_j\wedge v_i$
            \item $(a_iv_i+a_jv_j)\wedge (b_kv_k+b_lv_l)=(a_ib_k)v_i\wedge v_k+(a_ib_l)v_i\wedge v_l+(a_jb_k)v_j\wedge v_k+(a_jb_l)v_j\wedge v_l$
        \end{enumerate}
        And $A\wedge A$ where $A:V\to V$ is defined as 
        \begin{equation*}
            A\wedge A(v_i\wedge v_j)=Av_i\wedge Av_j
        \end{equation*}
        \item[(2)] Consider the matrix representation of $A=(A_{ij})$, on the basis $\{v_i\wedge v_j: i<j\}$, 
        \begin{align*}
            A\wedge A(v_i\wedge v_j)&=\sum_{k,l=1}^nA_{ki}A_{lj}(v_k\wedge v_l)\\
            &=\sum_{k<l}A_{ki}A_{lj}(v_k\wedge v_l)+\sum_{l<k}A_{ki}A_{lj}(v_k\wedge v_l)\\
            &=\sum_{k<l}A_{ki}A_{lj}(v_k\wedge v_l)-\sum_{l<k}A_{ki}A_{lj}(v_l\wedge v_k)
        \end{align*}
        Thus the diagonal term with respect to $v_i\wedge v_j$ is 
        \begin{equation*}
            A_{ii}A_{jj}-A_{ji}A_{ij}
        \end{equation*}
        Thus 
        \begin{equation*}
            \text{Tr}(A\wedge A)=\sum_{i<j}A_{ii}A_{jj}-A_{ji}A_{ij}
        \end{equation*}
        Now 
        \begin{equation*}
            \text{Tr}(A)^2=\sum_{i=1}^nA_{ii}^2+2\sum_{i<j}A_{ii}A_{jj}
        \end{equation*}
        and 
        \begin{align*}
            \text{Tr}(A^2)&=\sum_{k,l=1}^nA_{lk}A_{kl}\\
            &=\sum_{i=1}^nA_{ii}^2+2\sum_{k<l}A_{lk}A_{kl}
        \end{align*}
        Thus we see that 
        \[tr(\wedge^{2}A)=\frac{1}{2}(tr(A)^{2}-tr(A^{2}))\]
    \end{itemize}

\end{proof}

\begin{prob}
    Find a table of characters for the alternating group \(A_{5}\).
\end{prob}
\begin{proof}
    \[
\begin{array}{c|ccccc}
 & \text{1} & \text{20} & \text{15} & \text{12} & \text{12} \\ 
 & \text{Id} & (1\,2\,3) & (1\,2)(3\,4) & (1\,2\,3\,4\,5) & (1\,2\,3\,5\,4) \\ 
\hline
\chi_1 & 1 & 1 & 1 & 1 & 1 \\ 
\chi_2 & 3 & 0 & -1 & \phi & 1-\phi \\ 
\chi_3 & 3 & 0 & -1 & 1-\phi & \phi \\ 
\chi_4 & 4 & 1 & 0 & -1 & -1 \\ 
\chi_5 & 5 & -1 & 1 & 0 & 0 \\ 
\end{array}
\]
where $\phi=\frac{1+\sqrt{5}}{2}$.
\end{proof}




\end{document}